% Ad Familiares 7.10 --- headnote
% ad-familiares-07-10, ca. mid-December 54 BC, Rome

\textit{Cicero to C. Trebatius Testa, written from Rome late in
December 54 BC. Trebatius is wintering with Caesar's army in
Gaul, having been placed on the general's staff at the start of
the year by Cicero's recommendation (\textit{Fam.} 7.5). The
letter is one of the running pieces of the Trebatius sequence
and is almost entirely jest: a congratulation on Caesar's having
discovered that the young man is a real jurist, a fretful note
about the Gallic winter, and a parade of the standing in-jokes
of the cycle --- Britain (where gold and silver were rumoured and
turned out not to exist), Trebatius's refusal to swim in the
Ocean, his refusal to watch the British \textit{essedarii} or
chariot-fighters, his old appetite back home for even the
lowest of gladiator shows (the blindfolded \textit{andabatae}).
\textit{Mucius} and \textit{Manilius} are the great republican
jurists Q. Mucius Scaevola and M'. Manilius, invoked here to
sign off, deadpan, on the proposition that a cold lawyer should
sit by the fire.}

\textit{The third and fourth sections drop the joking and show
why Cicero kept writing these letters at all. He has been
prompting Caesar by post to advance Trebatius's interests, and
he now wants a report on what is working and what Trebatius
intends. The closing flourish --- one meeting between us, serious
or jocular, would be worth more than not just our enemies but
even our brothers the Aedui --- is at Caesar's expense: the Aedui
were the Gallic tribe Rome had formally enrolled as
\textit{fratres}, brothers of the Roman people, and Caesar's
\textit{De Bello Gallico} makes much of them. Cicero pretends to
adopt the official idiom and at the same time to find his real
brotherhood across the Alps with a lukewarm Trebatius rather
than with anyone in the camp.}
